\chapter{Resultados e discussão}

Este capítulo apresenta os resultados obtidos em relação aos requisitos e aos objetivos, discute as
decisões de projeto à luz desses resultados e expõe as limitações identificadas.

\section{Atendimento aos requisitos funcionais}

Os 42 requisitos funcionais especificados no Capítulo 4 foram implementados. Destes, 41 são
verificados por testes automatizados. O requisito de instalação (RF36) é verificado por um roteiro
próprio, que instala o pacote em uma pasta temporária e confere seu funcionamento sem alterar o
sistema. Algumas funções dependem de serviços externos ou de hardware, e sua verificação automatizada
usa dublês: a conversa livre com o modelo de linguagem (RF24), o reconhecimento remoto (RF03 e RF05) e
a voz neural (RF06). A escolha do microfone (RF39) foi conferida também no servidor de áudio real.
Nessas funções, a validação em uso real complementou os testes. A rastreabilidade
completa está no Apêndice A.

\section{Atendimento aos requisitos não funcionais}

A Tabela~\ref{tab:desempenho} compara as medições de desempenho da arquitetura inicial, serial, com as
da arquitetura em pipeline introduzida na versão 2.0.0 e mantida até a versão final.

\begin{singlespace}
\legendatabela{Medições de desempenho}\label{tab:desempenho}
\small
\begin{tabularx}{\textwidth}{LP{4.2cm}P{4.2cm}}
\toprule
\textbf{Medida} & \textbf{Arquitetura inicial} & \textbf{Arquitetura final} \\ \midrule
\endfirsthead
\toprule
\textbf{Medida} & \textbf{Arquitetura inicial} & \textbf{Arquitetura final} \\ \midrule
\endhead
Janela visível após abrir & 0,5 a 4 s & 0,8 s \\
Primeira palavra (modo offline, texto) & resposta e voz inteiras & cerca de 1,3 s \\
Primeira palavra (modelo de linguagem) & geração e síntese completas & primeira frase gerada e sintetizada \\
Tempo para interromper & apenas durante a reprodução & cerca de 10 ms, em qualquer estágio \\
Renderizar o avatar sem cache (crachá / \en{selfie}) & 2,9 s / 4,3 s & 1,8 s / 2,2 s, em segundo plano \\
Ler o avatar do cache & cerca de 0,4 s & 0,08 s / 0,2 s \\
Repetir uma fala já sintetizada & nova síntese & 0,1 ms (cache LRU) \\
Atraso de um lembrete & até 15 s & menos de 0,1 s \\
Pausa de fim de fala & 0,9 s fixos & 0,7 s, ignorando ruídos curtos \\
\bottomrule
\end{tabularx}
\fonte{elaboração própria (2026).}
\end{singlespace}

O Incremento 16 (versão 3.4.0) acrescentou um refinamento guiado por medições do aplicativo real,
apresentado na Tabela~\ref{tab:desempenho34}. O consumo em repouso foi medido em uma janela de 20 a 30 s
depois do fim da saudação falada, com as duas versões nas mesmas condições e com o cache do avatar já
pronto.

\begin{singlespace}
\legendatabela{Medições de desempenho e fluidez: versões 3.3.0 e 3.4.0}\label{tab:desempenho34}
\small
\begin{tabularx}{\textwidth}{LP{3.6cm}P{3.6cm}}
\toprule
\textbf{Medida} & \textbf{Versão 3.3.0} & \textbf{Versão 3.4.0} \\ \midrule
\endfirsthead
\toprule
\textbf{Medida} & \textbf{Versão 3.3.0} & \textbf{Versão 3.4.0} \\ \midrule
\endhead
Processamento em repouso & 29\% de um núcleo & 5\% de um núcleo \\
Custo de um quadro do palco (microbenchmark) & 13 ms & 5 ms \\
Janela visível após abrir & 1,7 a 2,4 s & 0,6 a 0,7 s \\
Primeira palavra da saudação & 3,1 a 3,2 s & 2,4 a 2,8 s \\
Duração da saudação & cerca de 10,1 s & cerca de 8,5 s \\
Silêncio entre frases faladas & cerca de 1,2 s & 0,15 s \\
Engasgo ao trocar de avatar & cerca de 170 ms & nenhum (no máximo 4 ms por quadro) \\
Memória residente & cerca de 369 MB & cerca de 368 MB \\
\bottomrule
\end{tabularx}
\fonte{elaboração própria (2026).}
\end{singlespace}

A memória não diminuiu: cerca de 150 MB correspondem às bibliotecas (o SDK do modelo de linguagem,
o reconhecedor local com seu modelo, o mixer de áudio e o processamento numérico) e cerca de 100 MB às
imagens já convertidas para a biblioteca gráfica. Um teste com a devolução explícita de memória ao
sistema operacional recuperou apenas 10 MB, o que indicou uso legítimo, e não fragmentação.

Os requisitos de desempenho (RNF01 a RNF05, RNF26 e RNF27) foram atendidos conforme as tabelas
anteriores. Os requisitos de confiabilidade (RNF06 a RNF09) foram atendidos e são verificados por testes que
simulam falhas de transação, adulterações do histórico, interrupções durante o fluxo de texto, recusas
e erros da API, ausência de microfone e falhas repetidas da escuta passiva. O requisito de exatidão da
ativação (RNF13) foi atendido no conjunto de calibração, com zero falsos positivos em 115 frases
parecidas, e confirmado em uso real com a voz do usuário depois da introdução do ganho automático.
Os requisitos de manutenibilidade (RNF19 a RNF22) foram atendidos e passaram a ser verificados
continuamente: nenhuma dependência fora da regra entre camadas, complexidade máxima de 10, tipos
verificados sem ocorrências, com todas as assinaturas tipadas, e cobertura de 85,9\%, acima do mínimo
de 80\%.

\section{Discussão}

\subsection{Latência percebida}

O maior ganho de experiência veio de uma decisão arquitetural, e não de otimizações locais. Ao
sobrepor a geração de texto, a síntese e a reprodução, o tempo até a primeira palavra deixou de
depender do tamanho da resposta. Esse resultado está de acordo com a observação de
\textcite{nielsen1993} de que respostas em torno de um segundo preservam o fluxo da interação. O custo
dessa escolha é a concorrência, que exigiu mecanismos explícitos: token de cancelamento por interação,
máquina de estados e barramento de eventos.

\subsection{Privacidade e exatidão da ativação}

O uso de um reconhecedor local com gramática restrita mostrou-se superior à transcrição livre para uma
frase de ativação curta e fora do vocabulário. As frases distratoras funcionam como uma classe de
rejeição explícita, o que explica a ausência de falsos positivos na calibração. Os testes com a voz
real, por outro lado, mostraram que a qualidade da captação é tão importante quanto o algoritmo: com o
microfone a 37\% de entrada, nenhum ajuste de gramática resolveria o problema. O ganho automático
melhorou a situação, mas não substitui um volume de entrada adequado.

\subsection{Memória auditável}

O encadeamento de resumos criptográficos, somado aos gatilhos de somente acréscimo, oferece duas
camadas de proteção: o próprio banco recusa alterações feitas pelo sistema e por ferramentas comuns, e a
verificação da cadeia detecta modificações feitas por fora, inclusive com o arquivo aberto em outro
programa. Como os registros guardam a descrição no momento do fato, a pergunta “o que eu fiz” é
respondida corretamente mesmo depois da remoção dos itens.

\subsection{Princípios SOLID na prática}

A revisão final mostrou que uma arquitetura em camadas bem definida não impede, por si só, a erosão
local da qualidade. Mesmo com a separação hexagonal desde a versão 2.0.0, as funções acumularam
condicionais, o orquestrador absorveu responsabilidades e duas dependências passaram a apontar no
sentido errado ao longo de nove incrementos. Os princípios SOLID funcionaram como critérios objetivos
para localizar e corrigir esses pontos, e as funções de aptidão de \textcite{ford2017} tornaram a
correção permanente. Em vez de depender da atenção de quem altera o código, a violação passa a ser
apontada pelo próprio processo de verificação.

Um efeito colateral positivo foi a melhoria da tipagem. Ao substituir acessos dinâmicos por contratos
explícitos, o verificador de tipos encontrou um caso real de uso possivelmente nulo, e a correção do
princípio de Liskov transformou um risco silencioso de perda de dados em um erro de tipo.

\subsection{Desenvolvimento assistido por inteligência artificial}

A experiência indica que um agente de programação acelera substancialmente a implementação, inclusive
de refatorações amplas, mas desloca a responsabilidade pela qualidade para as verificações automáticas
e para a validação com o usuário. Os defeitos mais difíceis do projeto foram identificados por medições
reais ou por execuções repetidas, como o volume baixo do microfone, a configuração que redirecionava a
API e o bloqueio da thread pelo coletor de lixo, e não pela leitura do código. Isso reforça o valor das
práticas de engenharia descritas neste trabalho, que se mostraram complementares ao uso do agente.

\section{Limitações}

As principais limitações identificadas são:
\begin{enumerate}[label=\alph*)]
  \item o reconhecimento de comandos e a voz neural dependem de serviços na internet; sem conexão, a
        conversa livre pode usar o modelo local, mas a voz de reserva também fica indisponível;
  \item o modelo local executado no processador leva de 15 a 45 s em pedidos que usam ferramentas e,
        por ser pequeno, só recebe ferramentas que criam ou consultam: concluir, remover e abrir
        programas exigem os comandos explícitos;
  \item a tela de configurações mostra os modelos disponíveis, mas a troca do modelo ainda é feita por
        variáveis de ambiente e vale ao reabrir o assistente;
  \item a restauração de um backup é manual, com o assistente fechado;
  \item a gramática restrita foi calibrada apenas para a frase “ok Fael”; outras frases usam a busca
        fonética livre, menos precisa;
  \item a calibração da ativação usou cinco vozes sintéticas e a voz de um único usuário real; uma
        avaliação com mais pessoas e ambientes ruidosos daria estimativas mais robustas de exatidão;
  \item os lembretes só são anunciados com o programa aberto; um lembrete vencido com o programa fechado
        é anunciado na abertura seguinte;
  \item o pacote de instalação foi construído e testado para Ubuntu 24.04 com Python 3.12; outras
        versões exigem reconstruí-lo na própria máquina, e as fotografias pessoais seguem dentro do
        pacote;
  \item as posições do rosto em cada fotografia foram medidas manualmente; um novo avatar exige
        cadastrar essas posições no catálogo.
\end{enumerate}
